\section{Models and Experimental Conditions}
\label{sec:models}

We evaluate seven Phase~2B configurations, named exactly as in the canonical experiment artifacts (\Cref{tab:models}).
Configurations differ by vendor model and, where applicable, thinking/mode switches exposed by the interface used for collection.
No configuration is treated as a gold standard.

\begin{table}[t]
\centering
\caption{Model configurations evaluated in all three regimes (canonical labels).}
\label{tab:models}
\begin{tabular}{@{}ll@{}}
\toprule
\textbf{Short label} & \textbf{Canonical configuration name} \\
\midrule
Gemini ON & Gemini 3.6 Flash --- extended thinking ON \\
Gemini OFF & Gemini 3.6 Flash --- extended thinking OFF \\
Claude Sonnet 5 & Claude Sonnet 5 --- Medium (default) \\
DeepSeek Expert ON & DeepSeek V3 Expert --- DeepThink ON \\
Qwen 3.8 Max Fast & Qwen 3.8 Max --- Fast \\
GPT-5.6 Luna & GPT-5.6 Luna --- thinking OFF \\
Grok 4.5 Fast & Grok 4.5 Fast \\
\bottomrule
\end{tabular}
\end{table}


\paragraph{Gemini ON vs.\ OFF.}
Two Gemini 3.6 Flash configurations differ only in the extended-thinking switch (ON vs.\ OFF).
Both use the multi-message Primed protocol described in \Cref{sec:repro}.

\paragraph{What is held constant.}
Within a regime, all configurations translate the same frozen Polish source under the same Direct/Primed definitions, with three independent repeats.
Cross-regime comparisons reuse the same configuration list.

\paragraph{What is not claimed.}
Vendor display names and mode switches are operational labels for reproducibility.
We do not equate interface labels with internal model checkpoints beyond what the operator interfaces exposed at collection time, and we do not rank configurations by quality.


\subsection{Collection interfaces as part of the method}
Because the experiment uses vendor chat/API surfaces rather than a single open-weights checkpoint under identical decoding controls, configuration identity includes the operator-facing mode switches recorded at collection time.
That is a feature of the ecological validity of the study and a limitation for mechanistic attribution to a specific open checkpoint.
The manuscript therefore reports configuration labels exactly and avoids rewriting them into unofficial synonyms.


\subsection{Non-goals of the configuration set}
The seven configurations were chosen as representative Phase~2B operator-facing setups, not as an exhaustive survey of all vendors or decoding hyperparameters.
Absence of an open-weights local baseline is a limitation for mechanistic replication, not an accidental omission in the Results tables.
